\section{Independent root review of the binary staircase theorem}
\label{r20pxroot:review}

\paragraph{Bound source and decision.}
This review binds
\path{inertia/product_extension_probe.tex}, SHA256
\texttt{b618a9e134f13e78b19f1e8a90c7a566e2f87f6138caf93765bbf3500accee91}.
The decision is \textbf{PASS relative to the explicitly cited classical
and frozen interfaces}.  The root reviewer reconstructed the geometry
and correspondence argument separately from the author's source and
read the historical R9 ST/TR/DB material as candidate evidence only.
The author and reviewers communicated; this is an independent internal
check, not external or blind peer review.  The historical idea and the
two-node calculation are not advertised as new discoveries.

\subsection{Geometry checked independently}

For the local product cone, its lattice has rank $r+s-1$.  A first
path ray together with the successive differences of path rays is a
basis: the differences give the two root lattices of vectors of sum
zero, and the first ray has common sum one.  Thus each path cone is
unimodular and all original rays have height one.  The mass transfer
calculation with heights $-ij$ strictly lowers the height on every
incomparable pair.  The unique monotone coupling with prescribed
marginals gives exactly the path simplices.  Global component orders
make the heights and subdivisions agree on faces.  This supplies a
projective regular reduced semistable modification, including all
surface triple strata.

The quotient star at a component ray splits into the lower and upper
orthants.  In each orthant, reverse-product star subdivisions at
$u_a+v_b$ give the path fan by induction on the final step.  A
different two-face subdivision leaves the unsplit two-face available.
Consequently these are actual smooth codimension-two boundary
blowups.  Their centers map birationally to the original products of
two double curves.  The surface birational-morphism factorization
into point blowups gives their abelian-type Chow motives.  The
ordinary blowup decomposition gives actual injection and retraction
maps and base projector $f^*f_*$.

On a same-first-label double stratum, all other first-factor
parameters have valuation zero on both rays.  Thus none can vanish
at the generic point of that double stratum.  There is no extra
component over the deleted first-factor boundary.  The map to the
original same-label double is proper birational of degree one.
When that label changes, the map instead factors through its actual
boundary.  These facts justify using the real double strata, rather
than a formal tensor product of two incidence matrices.

\subsection{Correspondences checked independently}

The transport lemma is valid on every input.  For
$\beta=f^*f_*\beta+\beta_0$, the exceptional term has
$f_*\beta_0=0$ and vanishes away from the modification locus.
Localization puts it in supported cohomology.  Restriction to the
double and pushforward to the base preserve this support, and the
transcendental surface projector annihilates it.  This proves the
restriction identity also on exceptional inputs.  Self-adjointness
and Poincare adjoints give the Gysin identity.  Neither flatness of
the birational square nor a Chow-level annihilation theorem was used.

For a surface, the only degree-two classes supported on a union of
boundary curves are its divisor classes.  The transcendental
projector kills these and restricts to zero on every boundary curve.
It is therefore legitimate to apply the lemma factorwise.
The actual Chow projectors $P^S$ and $P^T$ commute, because they
come from the same base product and $f_*f^*=1$.  The actual incidence
operator commutes with them on full cohomology.  On the double
transcendental sector it vanishes; on a single labelled sector it is
the identity on that coefficient tensored with the other surface's
actual incidence.  The diagonal follows from the same divisor
calculation and the principal total fiber, not from a tensor-sum
ansatz.

The residual $Q$ corner is an actual Chow summand: the base part is
a tensor product of the nontranscendental surface summands and the
exceptional part consists of the displayed abelian-type centers.
The final source explicitly makes a finite unramified extension
splitting the finite Neron--Severi and point Artin actions before
describing these as Tate pieces.  Thus a single finite abelian
carrier suffices for all entries and all needed twists.

The root reviewer checked Deligne--Clozel's original 2009 text,
Theorem~1 and Remark~(b), printed p.~496: one suitable prime works
for all powers of that one carrier.  Jannsen's original 1992
Theorem~1 is used only for numerical motives.  Numerical zero over
the native finite field remains zero after finite constants, by
norms and the full extension degree.  Hence numerical errors on
the residual carrier have zero full geometric class at the selected
prime.  The argument does not intersect unrelated prime sets and
does not assert an all-prime comparison theorem for arbitrary
abelian-type motives.

\subsection{Actual cycles and final scope}

The sum of the sector inverses is an actual rational
degree-minus-one correspondence with $A^H G^H A^H=A^H$.
Tensoring it with a good coefficient acts on all mixed cohomology.
For an actual algebraic lift $x$ of a numerical $A$-kernel class,
$x-G^HA^Hx$ remains actual algebraic, retains the numerical class,
and is in the cohomological $A$-kernel.  The algebraicity of the
preimage is explicit.  No algebraicity of all component cohomology
is assumed.

The odd weight--monodromy input follows from the known surface
case and the good factor by the tensor rule.  Applying the frozen
AK19/OW19 and R16 dominant detection, then full-degree finite
descent, gives the original-cycle obstruction vanishing.  VT17
supplies the ordinary-actual form with the exact original support
and all compact actual tests.  The proof reapplies to every newly
specified cycle after a finite valued extension.  There is no
assumption that an alteration is finite on the support complement.

The two-surface fourfold itself is already in SP19.  The present
proof's stated extension is the arbitrary potentially-good
coefficient with all mixed positive codimensions.  The R9 family
idea is acknowledged.  This review alone does not certify the
three-or-more-bad-surface induction, a model-independent $U_2=0$
statement, general strong G, or the full height domain for arbitrary
varieties.
